We embed this environment in a global-games framework in which agents receive heterogeneous private signals and observe a common public signal. Agents choose whether to acquire foreign currency for hedging and for speculative exit. The equilibrium is characterized by a cutoff strategy: there exists a threshold belief such that agents attack when their posterior exceeds this cutoff.

Improved public information facilitates coordination among agents. As the precision of the stablecoin-based signal increases, agents place more weight on the common signal and less on their idiosyncratic private information. This strengthens strategic complementarities and raises the likelihood that agents coordinate on a speculative attack.

Our central result is that stablecoins have a non-monotonic effect on welfare. At low levels of exchange rate misalignment, stablecoins improve welfare by reducing transaction costs and improving risk-sharing. However, as misalignment increases, the improved information and coordination facilitated by stablecoins increase the probability of a crisis, reducing welfare.

Quantitatively, the model delivers three key predictions. First, stablecoin demand increases with exchange rate overvaluation. Second, the presence of stablecoins increases crisis probabilities relative to a cash-only benchmark. Third, welfare exhibits a non-monotonic relationship with exchange rate misalignment.

Our paper contributes to three strands of literature. First, it relates to the literature on currency crises and global games, where public information can facilitate coordination and destabilize equilibria. Unlike standard models, the precision of the public signal in our framework is endogenous. Second, we contribute to the emerging literature on digital currencies by emphasizing their role as information and coordination devices rather than purely transactional technologies. Third, we connect to the literature on parallel exchange rate markets by introducing a new instrument that alters both market microstructure and information flows.

More broadly, our results suggest that stablecoins can transform parallel exchange rate markets into transparent, high-frequency coordination environments. While they improve efficiency, they may also increase macro-financial fragility by facilitating coordinated exits.
